\section[Rédaction Bourbaki \redac{402} : distances localisables]{Rédaction Bourbaki \redac{402} : distances localisables et chaînes\protect\verifie{Redaction402}}\label{sec:bourbaki-402}

La rédaction \redac{402}, \q{Géométrie différentielle. Chapitre~II : Géodésiques}, des archives Bourbaki (Archives Henri Poincaré, \ahp{953}) est datée de juin 1964 par sa notice, qui la donne à Jean-Louis Koszul ; elle compte 127 pages. Son §~5 s'ouvre (p.~10--11 du tapuscrit) sur une notion de pure topologie métrique, qui servira à dire que la distance d'une variété riemannienne se calcule localement. Sa transcription reste locale pour l'instant ; on n'en cite ici aucun passage.

\emph{Chaînes et longueur.} Soient $X$ un espace muni d'une distance $d$ et $W$ un voisinage de la diagonale de $X \times X$. Une $W$-chaîne de $a$ à $b$ est une suite finie $a = x_0, x_1, \dots, x_n = b$ dont chaque pas $(x_i, x_{i+1})$ est dans $W$ ; sa longueur est $\sum_i d(x_i, x_{i+1})$. Par l'inégalité triangulaire, une chaîne n'est jamais plus courte que $d(a, b)$. La définition~1 dit que $d$ est \emph{localisable} si, pour tout voisinage $W$ de la diagonale et tous $a$, $b$, la distance $d(a, b)$ est la borne inférieure des longueurs des $W$-chaînes de $a$ à $b$.

\emph{Un exemple.} Sur le cercle unité du plan, la distance euclidienne (la corde) n'est pas localisable. Une chaîne à petits pas suit le cercle, et la somme de ses cordes est voisine de la longueur d'arc parcourue, qui dépasse la corde de $a$ à $b$ dès que $a \neq b$. La distance géodésique, longueur du plus petit arc, est au contraire localisable : les chaînes qui suivent l'arc ont exactement sa longueur. Être localisable, c'est donc dire que la distance se mesure en mettant bout à bout de petits pas, comme une longueur de chemin. Cet exemple n'est pas formalisé.

\emph{Le choix de la définition.} Le fichier écrit la borne inférieure sous la forme suivante : pour tout $\delta > 0$, il existe une $W$-chaîne de longueur $< d(a, b) + \delta$. Cette forme garde l'existence des chaînes, qu'une borne inférieure sur un ensemble vide de réels cacherait. La rédaction quantifie sur \emph{tous} les voisinages de la diagonale, et non sur les seuls voisinages uniformes $\{d < \varepsilon\}$ ; la démonstration de la proposition~2 a besoin d'un voisinage qui n'est pas uniforme. Une note manuscrite en regard propose la forme en $\varepsilon$. Les deux ne sont pas équivalentes, comme on le verra.

\begin{encadre}
\begin{enonce}[Bourbaki, rédaction \redac{402}, §~5, propositions~1 et~2, p.~11]
Soit $X$ un espace pseudo-métrique dont la distance est localisable.
\begin{enumerate}
\item[a)] (Proposition~1.) Une seconde distance localisable, qui définit la même topologie et coïncide avec $d$ sur un voisinage de la diagonale, est égale à $d$.
\item[b)] (Proposition~2.) Si $a \in A$ et $b \notin A$, alors $d(a, b)$ est la borne inférieure de $d(a, c) + d(c, b)$ pour $c$ dans la frontière de $A$.
\item[c)] (Corollaires.) On a $d(a, F) + d(F, b) \leq d(a, b)$, où $F$ est la frontière de $A$. La frontière de la boule ouverte de centre $a$ et de rayon $r > 0$ est la sphère, et son adhérence est la boule fermée. Si la boule fermée de rayon $r$ est compacte et $d(a, b) \geq r$, il existe $c$ avec $d(a, c) = r$ et $d(c, b) = d(a, b) - r$.
\end{enumerate}
L'espace est alors connexe. La rédaction le suppose connexe ; avec la définition par tous les voisinages, cette hypothèse est une conséquence et n'est pas nécessaire.
\end{enonce}
\end{encadre}

\begin{proof}
a) Soient $d_1$ et $d_2$ les deux distances, égales sur le voisinage $W$. Pour une $W$-chaîne de $a$ à $b$, on a $d_1(a, b) \leq \sum d_1(x_i, x_{i+1}) = \sum d_2(x_i, x_{i+1})$, la première inégalité par l'inégalité triangulaire, l'égalité parce que chaque pas est dans $W$. Comme $d_2$ est localisable, le membre de droite peut être rendu $< d_2(a, b) + \delta$ pour tout $\delta > 0$. Donc $d_1 \leq d_2$, et par symétrie $d_1 = d_2$.

b) Le cœur est le suivant : pour tout $\varepsilon > 0$, il existe $c$ dans la frontière avec $d(a, c) + d(c, b) < d(a, b) + \varepsilon$. Si $a$ ou $b$ est sur la frontière, on le prend pour $c$. Sinon, $a$ est intérieur à $A$ et $b$ extérieur. Soit $V$ l'ensemble des points à distance $< \varepsilon/3$ de la frontière, et $W$ la réunion de trois carrés : intérieur de $A$ fois lui-même, $V \times V$, extérieur de $A$ fois lui-même. C'est un ouvert, et il contient la diagonale, puisque tout point est intérieur, sur la frontière ou extérieur. Une $W$-chaîne de longueur $< d(a, b) + \varepsilon/3$ part de l'intérieur et finit à l'extérieur. Un pas dans le premier carré reste à l'intérieur, et un pas ne peut passer de l'intérieur à l'extérieur que par $V \times V$ ; la chaîne a donc un point $x_i$ dans $V$, à distance $< \varepsilon/3$ d'un point $c$ de la frontière. En coupant la chaîne en $x_i$, on a $d(a, x_i) + d(x_i, b) < d(a, b) + \varepsilon/3$, et deux inégalités triangulaires par $c$ donnent la conclusion. L'inégalité inverse, $d(a, b) \leq d(a, c) + d(c, b)$, est l'inégalité triangulaire.

c) Les corollaires en découlent. Pour la boule ouverte, un point $b$ de la sphère est approché par des points $c$ de la frontière de la boule ; une telle frontière est toujours contenue dans la sphère. Pour le dernier, la frontière de la boule est un fermé du compact, et la fonction $c \mapsto d(c, b)$ y atteint son minimum, qui vaut $d(a, b) - r$ par b). Enfin, si $X$ était réunion de deux ouverts fermés disjoints non vides $A$ et son complémentaire, b) donnerait un point de la frontière de $A$, qui est vide : $X$ est connexe.
\end{proof}

\emph{La note manuscrite.} Appelons $\varepsilon$-localisable une distance qui vérifie la définition pour les seuls voisinages $\{d < \varepsilon\}$. Sur $\mathbf Q$, la distance usuelle est $\varepsilon$-localisable : une progression arithmétique de $a$ à $b$ en $n$ pas a pour longueur exactement $|b - a|$, avec des pas aussi petits qu'on veut. Mais $\mathbf Q$ n'est pas connexe, donc n'est pas localisable, et la proposition~2 y est fausse : pour $A = \{x < \sqrt 2\}$, ouvert et fermé dans $\mathbf Q$, la frontière est vide. La forme proposée en regard ne suffit donc pas à la démonstration de la page.

\begin{sloppypar}
\noindent\emph{Vérification en Lean.} La définition est \texttt{IsLocalizable} (pour une fonction $d$ sur un espace topologique) et \texttt{Localizable} ; la proposition~1 est \texttt{eq\_of\_isLocalizable} et \texttt{PseudoMetricSpace.\allowbreak{}eq\_of\_localizable} ; la proposition~2 est \texttt{Localizable.\allowbreak{}exists\_mem\_frontier} et \texttt{Localizable.\allowbreak{}dist\_eq\_iInf\_frontier} ; les corollaires sont \texttt{infDist\_add\_infDist\_le}, \texttt{frontier\_ball}, \texttt{closure\_ball} et \texttt{exists\_dist\_eq\_sub} ; la connexité est \texttt{Localizable.\allowbreak{}preconnectedSpace}, et l'exemple de $\mathbf Q$ est \texttt{epsLocalizable\_rat}, \texttt{not\_localizable\_rat} et \texttt{not\_exists\_mem\_frontier\_rat}, dans \href{\leanblobmain/Redaction402.lean}{\texttt{Redaction402.\allowbreak{}lean}} ; Lean y tire l'échec de la proposition~2 de la non-connexité de $\mathbf Q$, sans exhiber la partie $A$. Pour la proposition~1, Lean n'utilise de $d_1$ et $d_2$ que l'inégalité du polygone, et non leurs autres propriétés de distance.
\end{sloppypar}

Ne sont pas démontrés ici : l'usage du §~5 dans la suite du chapitre, à savoir que la distance d'une variété riemannienne est localisable ; l'exemple du cercle.

\litterature{le relevé des rédactions la classe parmi les textes retrouvés. Les distances localisables au sens des voisinages sont proches des distances de longueur (ou intrinsèques), que Burago, Burago et Ivanov caractérisent au §~2.4 de leur cours par l'existence de milieux approchés \cite{bbi2001} ; la forme par chaînes, pour des espaces non complets, n'a pas été comparée à leur énoncé. Mathlib n'a pas de théorie des espaces de longueur.}
